\documentclass[letterpaper]{article}
\usepackage{iccc}
\usepackage{times}
\usepackage{helvet}
\usepackage{courier}
\usepackage{booktabs}

\pdfinfo{
/Title (Exploratory Creativity in Literary Interpretation)
/Subject (Proceedings of ICCC)
/Author (Anonymous)}

\title{Exploratory Creativity in Literary Interpretation:\\
Manipulating Critical Lenses at Zero Marginal Cost}
\author{Anonymous}
\setcounter{secnumdepth}{0}

\begin{document}
\maketitle
\begin{abstract}
\begin{quote}
We present a system for generating multi-voice literary discussions
about novels that decomposes interpretation into independently
manipulable components: expert personas (critical frameworks),
passage selection (reading strategies), and conversational
scaffolding.  Testing across 15 novels and 180 conditions, we find a
\emph{convergence paradox}: two algorithmically distinct reading
strategies---one synoptic, one analytical---select nearly disjoint
passage sets (Jaccard~0.057) yet produce discussions whose character
is determined by the interpretive framework, not the reading.  This
makes the system a tool for \emph{exploratory creativity} in Boden's
sense: systematic traversal of an interpretive space defined by
critical perspectives and textual affordances, at zero marginal cost
per exploration.  Removing textual grounding reveals that ungrounded
generation produces not random output but graduated \emph{pastiche}
that recombines the novel's vocabulary---a form of combinational
creativity whose fidelity tracks training-data exposure.  We argue
that the system's value lies not in producing creative artefacts but
in making the \emph{space of possible interpretations} navigable.
\end{quote}
\end{abstract}

\section{Introduction}

Computational creativity research has long distinguished between
\emph{exploratory} creativity---traversing a defined conceptual
space---and \emph{combinational} creativity---producing novel
combinations of familiar elements \cite{boden2004creative}.  Most
LLM-based creative systems operate in combinational mode: they
generate novel text by recombining learned patterns.  We present a
system that enables exploratory creativity in a specific domain:
literary interpretation.

The system generates structured multi-voice podcast discussions about
novels.  Three expert personas, each embodying a distinct critical
framework (e.g.\ formalist, social-historical, reader-response),
discuss passages selected by an algorithmic reading strategy.  The key
insight is that the components are independently variable: the same
novel can be discussed through different lenses, using different
passage selections, with different conversational scaffolding---and
each variation is a distinct point in an interpretive space that a
human editor can explore.

This exploration is \emph{cheap}.  The passage-selection step uses
min-cost network flow \cite{ahuja1993network} and runs in under 100ms
with no LLM calls.  An editor can preview dozens of interpretive
configurations before committing to the expensive generation step.
This ``preview-before-commit'' workflow distinguishes the system from
generative tools like NotebookLM \cite{notebooklm2024}, which produce
a single output with no mechanism for systematic variation.  It
realises the kind of ``shared collaborative space'' advocated by
\citeauthor{ocampo2025beyond}~\shortcite{ocampo2025beyond}, in which
a human editor has genuine agency over the creative parameters rather
than merely prompting a chat interface.

\section{The Convergence Paradox}

Our central empirical finding concerns the relative power of different
creative levers.  We implement two reading strategies as
passage-selection algorithms applied to 15 Victorian novels (60,774
passages, 7 authors, 1850--1924):

\begin{itemize}
\item \textbf{Synoptic}: min-cost flow selects 59 passages from 28
chapters, favouring character development (+8.3~pp) and plot (+10.2~pp).
\item \textbf{Analytical}: contextual embeddings select 32 passages
from 20 chapters, favouring social critique (+13.4~pp) and thematic
depth (+6.9~pp).
\end{itemize}

These are recognisably different scholarly strategies---the seminar
leader who assigns the whole novel versus the one who curates a dossier
of key passages.  They agree on fewer than 6\% of their selections
(mean Jaccard~0.057, range 0.000--0.118 across 28 novel$\times$panel
pairs).  Yet the resulting discussions converge: expert airtime is
stable to $\pm$2\% across all 15 novels and both strategies.  Eleanor
Hartley (literary critic) receives 32--37\% of expert words, James
Blackstone (social historian) 30--37\%, Caroline Woodcourt (close
reader) 29--34\%, regardless of which algorithm selected the passages
or which novel is being discussed.

This stability was not predicted.  We expected expert prominence to
\emph{adapt} to the novel's affordances---that a novel rich in social
critique would give the social historian more airtime.  The data shows
the opposite: frameworks are robust, maintaining their character across
15 novels spanning Dickens, Eliot, Gaskell, Forster, Collins, Gissing,
and Oliphant.

In Boden's framework, this means the \emph{conceptual space} of
literary discussion is shaped primarily by the interpretive
framework---the critical lens---not by the input material.  The space
is framework-defined, not text-defined.  This is a structural finding
about where creative agency resides in the system: it resides in the
choice of perspective, not in the selection of evidence.

\section{Pastiche as Combinational Creativity}

What happens when textual grounding is removed entirely?  In our
\emph{no-passages} condition, the LLM must discuss the novel from
prior knowledge alone.  The result is not silence or incoherence but
graduated \emph{pastiche}: the system constructs plausible quotations
by recombining the novel's own vocabulary.  Of 726 invented quotes,
82\% draw 90--100\% of their individual words from the novel's word
stock.  The mean overlap is 95.1\%.

This is combinational creativity in a precise sense: the system
produces novel combinations of familiar elements (the novel's
vocabulary) that are locally plausible but globally unfaithful.

\textbf{The knowledge gradient.}  The fidelity of pastiche tracks
cultural prominence: \emph{Middlemarch} (57\% verified),
\emph{Bleak House} (53\%), \emph{David Copperfield} (44\%) at the
top; \emph{New Grub Street} (20\%), \emph{The Odd Women} (6\%),
\emph{Hester} (0\%) at the bottom.  All 15 novels are on Project
Gutenberg and plausibly in the LLM's training data; the gradient
reflects not textual presence but the volume of secondary discussion
(study guides, academic papers, Wikipedia) that reinforces specific
passages \cite{carlini2021extracting,carlini2023quantifying}.
\citeauthor{bayard2007comment}~\shortcite{bayard2007comment} argued
that humans navigate literary culture through graduated approximation;
the LLM does something analogously measurable, though by entirely
different mechanisms.

\begin{table}[t]
\centering
\small
\caption{The knowledge gradient: ungrounded quote verification (\%)
by novel.  All novels are on Project Gutenberg; the gradient reflects
analytical reinforcement, not textual absence.  Grounded conditions
achieve 97--98\% for all novels.}
\label{tab:gradient}
\begin{tabular}{llr@{\hspace{6pt}}r}
\toprule
Novel & Author & Nop\% & Gap \\
\midrule
Middlemarch & Eliot & 57 & 42\,pp \\
Bleak House & Dickens & 53 & 45\,pp \\
David Copperfield & Dickens & 44 & 53\,pp \\
Hard Times & Dickens & 39 & 59\,pp \\
Passage to India & Forster & 39 & 55\,pp \\
Mill on the Floss & Eliot & 37 & 59\,pp \\
Daniel Deronda & Eliot & 30 & 69\,pp \\
Our Mutual Friend & Dickens & 29 & 69\,pp \\
New Grub Street & Gissing & 20 & 79\,pp \\
The Odd Women & Gissing & 6 & 91\,pp \\
Hester & Oliphant & 0 & 98\,pp \\
No Name & Collins & 0 & 95\,pp \\
\bottomrule
\end{tabular}
\end{table}

\textbf{Within-author evidence.}  The gradient persists within
authors, controlling for prose style.  Among Dickens's four novels:
\emph{Bleak House} (53\%) $>$ \emph{David Copperfield} (44\%) $>$
\emph{Hard Times} (39\%) $>$ \emph{Our Mutual Friend} (29\%).  Among
Eliot's three: \emph{Middlemarch} (57\%) $>$ \emph{Mill on the Floss}
(37\%) $>$ \emph{Daniel Deronda} (30\%).  Same author, same prose
style, different cultural prominence, different pastiche fidelity.

For computational creativity, this finding matters because it
characterises a \emph{failure mode} as a \emph{creative mode}.  The
system is not hallucinating randomly; it is performing pastiche---imitation
of an author's manner without direct quotation.  This is a form of
creativity: it produces novel, plausible, stylistically coherent text.
But it is creativity unmoored from its source, and the graduated nature
of its infidelity makes it measurable through our five-tier scheme
(verified, paraphrase, distant echo, no clear source, invented).

\section{Exploring Interpretive Space}

The system's practical value for computational creativity lies in its
support for exploratory traversal of interpretive space.  An editor
can:

\begin{enumerate}
\item \textbf{Swap critical lenses}: replace a formalist with a
Marxist, an astronomer, or an NLP scientist, and observe how the
discussion reshifts.

\item \textbf{Adjust demand profiles}: increase one expert's appetite
for social critique or narrative technique, producing graduated
passage-set changes (Jaccard 0.5--0.8 with baseline).

\item \textbf{Add or remove scaffolding}: host preparation raises
question density from $\sim$1 to $\sim$11 per segment, uniformly
across all novels.

\item \textbf{Compare grounded and ungrounded}: observe the pastiche
gradient as grounding is removed, measuring how much of the
discussion's substance depends on actual textual evidence.
\end{enumerate}

Each exploration incurs zero marginal LLM cost at the
passage-selection level.  The expensive step (script generation) is
deferred until the editor has identified a promising configuration.
This is Boden's \emph{exploratory creativity} operationalised as a
human-AI workflow: the human defines the space (through persona and
demand choices), the machine traverses it (through optimisation and
generation), and the human evaluates the results.

\section{Relation to Computational Creativity}

We see three connections to the CC literature:

\textbf{Creativity resides in the framework, not the output.}
\citeauthor{colton2012computational}~\shortcite{colton2012computational}
argue that attributing creativity requires understanding the
\emph{process}, not just evaluating the \emph{product}.
\citeauthor{morain2025prompt}~\shortcite{morain2025prompt} ask
whether prompt engineering is the ``creativity knob'' for LLMs.  Our
convergence paradox provides an empirical answer: the persona
prompt---not the textual input or retrieval method---determines the
character of the output.  The creative act is the choice of lens.

\textbf{Pastiche is a measurable creative mode.}
\citeauthor{ritchie2007}~\shortcite{ritchie2007} proposes criteria
for attributing creativity to programs, including novelty and value.
Pastiche is novel (the sequences do not exist in the source text) and
locally valuable (they are stylistically coherent).  But it lacks
\emph{fidelity}---it does not faithfully represent the source.  Our
five-tier verification scheme (verified, paraphrase, distant echo, no
clear source, invented) provides a graded measure of this trade-off.

\textbf{The system enables human exploratory creativity.}
\citeauthor{ismayilzada2025evaluating}~\shortcite{ismayilzada2025evaluating}
find that LLMs generate stylistically complex stories but fall short
in novelty and diversity compared to humans.
\citeauthor{doshi2024generative}~\shortcite{doshi2024generative} find
that generative AI enhances individual creativity but reduces
collective diversity.  Our system addresses both concerns: by making
the interpretive framework an explicit, manipulable parameter, it
produces \emph{diverse} outputs from the same source material.  The
system does not converge on a single ``best'' interpretation; it maps
the space of possible interpretations.

\section{What If I Put My Friends in the Podcast?}

The system's personas are not limited to literary scholars.  To test
the boundaries of exploratory creativity, we replaced the default
panel (a formalist, a social historian, and a close reader) with three
American professionals from outside the humanities: a computer
scientist specialising in speech recognition and NLP, an
observational astronomer who studies stellar formation, and a
musicologist who studies Cold War cultural diplomacy.

The resulting discussions of \emph{Bleak House} are recognisably
different from the literary-critical baseline.  The astronomer reads
Dickens's fog as a physical formation condition---opacity is not
metaphor but the baseline state of complex systems.  The NLP
scientist analyses agent deletion in the opening sentences: ``There is
no main clause.  It is a sequence of noun phrases and participial
clauses with no finite verb until paragraph two.''  The musicologist
compares Chancery to Soviet institutions that consumed the artists
they were supposed to support.

Two observations matter for computational creativity:

First, the \emph{convergence paradox still holds}.  Despite radically
different personas, expert airtime remains balanced and the
discussions are structurally similar to the literary-critical ones.
The framework dominates even when the framework is non-literary.

Second, the experts bring their \emph{metaphors} but not their
\emph{methods}.  The astronomer does not propose measurements or
hypotheses; she produces literary-critical discourse decorated with
astronomical imagery.  This is a limitation of the prompt
architecture: the host preparation questions (``what strikes you
about these passages?'') import literary-critical norms regardless
of the persona.
\citeauthor{sood2025conversational}~\shortcite{sood2025conversational}
ask whether conversational interfaces limit creativity; our finding
suggests that conversational \emph{scaffolding} can constrain
disciplinary expression even when the persona is unconstrained.
A prompt designed for genuinely interdisciplinary
encounter---``what would your discipline's tools reveal about this
passage?''---might produce different results.  This is itself a
finding about where creative constraint operates: not just in the
persona but in the conversational scaffold.

\section{Debate Spark: Is Systematic Variation Creative?}

The convergence paradox makes one thing obvious: the personas and
framework bear the creative load in this system---perhaps badly.  The
LLM is a \emph{realiser}.  It takes a specification (persona, passages,
scaffolding instructions) and renders it into fluent multi-voice
discourse.  The rendering is impressive.  But the creative decisions---which
lens to apply, which dimensions to weight, which experts to put in
the room together---are made by the human editor before the LLM is
invoked.

This is not a controversial claim.  It is the direct implication of
the convergence paradox: since the interpretive framework determines
the output's character regardless of reading strategy, the framework
is where the creative agency sits, and the framework is
human-authored.

What \emph{is} debatable is whether the system's ability to
\textbf{vary systematically at scale} constitutes a creative
contribution that deserves the CC community's attention.  We think it
does, for three reasons:

\textbf{1.~Scale makes comparison possible.}  A human editor cannot
commission 180 literary discussions to compare how different lenses
reveal different aspects of 15 novels.  The system can, and the
comparison itself---the ability to see that a Marxist reading and a
formalist reading of \emph{Hard Times} select completely different
passages but produce structurally similar discussions---is a finding
that no single discussion could produce.  Creativity here lies not in
any individual output but in the \emph{landscape of outputs}.

\textbf{2.~Zero-cost exploration enables serendipity.}  Because
passage selection runs in under 100ms with no LLM calls, an editor
can try configurations they would never commission at full cost.
``What if I put an astronomer in the panel?'' is a question you ask
when asking is free.  The interdisciplinary panel (Section~7) emerged
from exactly this dynamic: a whimsical configuration that produced
genuinely surprising readings.

\textbf{3.~The pastiche gradient reveals something new.}  The
graduated degradation from faithful quotation (57\% for
\emph{Middlemarch}) to pure invention (0\% for \emph{Hester}) is not
something the system was designed to produce.  It is an emergent
property of removing grounding from an LLM that has differentially
memorised its training data.  Whether this counts as ``creativity'' is
precisely the kind of boundary question CC research should address.

\textbf{Objections we take seriously:}

The strongest objection is that systematic variation is
\emph{engineering}, not creativity.  A tool that renders many versions
of the same specification is an industrial process, not an artistic
one.  We would accept this objection if the variations were
predictable---but they are not.  The convergence paradox was
surprising.  The pastiche gradient was surprising.  The
interdisciplinary panel's ``metaphors but not methods'' finding was
surprising.  The system produces surprises not through generative
novelty but through \emph{systematic comparison}---and we argue that
surprise-through-comparison deserves a place in the CC vocabulary
alongside surprise-through-generation.

A second objection is that calling the human the ``creative agent''
diminishes the LLM's contribution.  The LLM produces the actual
prose---hundreds of thousands of words of fluent, contextually
appropriate literary discussion.  Is that not creative?  We think not,
for the reason the convergence paradox demonstrates: the prose is
\emph{determined} by the specification.  Change the persona and the
prose changes predictably.  A system whose outputs are
specification-determined is a realiser, however impressive its
realisations.  But we recognise this is a contested claim and invite
the community's response.

\section{Limitations}

Our ``interpretive frameworks'' are LLM persona prompts, not cognitive
models.  As \citeauthor{shanahan2024talking}~\shortcite{shanahan2024talking}
warns, we should not map LLM behaviour onto human cognitive categories
without justification.  The personas are stereotypes, not
theories---useful experimental instruments, not models of creative
cognition.  We report automated metrics only, with no human evaluation
of creative quality.

\section{Author Contributions}
Anonymised for review.

\bibliographystyle{iccc}
\bibliography{iccc_refs}

\end{document}
